% P54 | 6 Comparison of methods | Which method for which job
% guide: dev_notes/paragraph_guides/P54_which_method_for_which_job.txt
For a vanilla European option the closed form is the answer when it exists, and otherwise a Leisen--Reimer tree or a smoothed and extrapolated CRR tree is faster and more accurate than any simulation.
Monte Carlo serves there only as an independent cross-check.
For an American option on one underlying a CRR tree with Richardson extrapolation delivers the price, the exercise boundary and the hedge in tens of milliseconds, and a PDE solver is the better choice when high accuracy is needed.
The Barone--Adesi--Whaley formula is good enough for screening and for initialising a root-finder.
For path-dependent payoffs or several underlyings, Monte Carlo with variance reduction is the default, with a regression step such as Longstaff--Schwartz if early exercise is allowed.
The cost of a tree or a grid grows exponentially with the dimension and rules them out beyond two or three underlyings.
Sensitivities come from the same lattice or grid as the price where possible, and from pathwise or likelihood-ratio estimators in Monte Carlo.
A stress of the volatility input is more informative than any refinement of the numerics, because in both parts of this report the effect of a wrong volatility dwarfs the numerical error of a reasonable method.
